\chapter{活细胞做的计算机}
% 完整版：chapters/21_living.tex

\pencilfig{21}{\pencilart{21}}{芯片上的活神经元：下面有成千上万个电极，每一个都既能听也能说。}

能不能不用硅，而用活的神经元来造机器？能，而且已经有人在造了。神经元直接长在带有成千上万个电极的芯片上——长成平平的一层，或者长成一个个像微型大脑的小团块。每个电极都既能听，也能送出信号。这是一个电路完全敞开的试验台：一切都能测量，一切都能干预。只是里面没有电台——它是一条没有控制寄存器的数据总线。

\section*{网络自己会做什么}

放任不管的时候，培养物会时不时整体爆发一次：几乎所有神经元一起咔嗒，然后安静下来——间隔从一秒到几分钟不等。这是我们熟悉的振荡器：网络自我兴奋，然后疲劳，直到连接补足物质的储备。

\pencilfig{129}{%
% spike raster of 8 neurons in a culture left to itself: sparse single clicks and shared bursts
\foreach \b in {0.9,3.1,4.0,6.6}{\fill[pcAmber, opacity=0.18] (\b-0.1,-0.15) rectangle (\b+0.32,2.95);}
\foreach \y/\ss in {0/{0.96,1.0,1.13,2.43,3.0,3.12,3.24,4.02,4.09,4.15,6.66,6.69,6.81},
  1/{0.46,0.88,1.0,1.05,3.09,3.2,3.94,4.04,4.37,6.65,6.71},
  2/{0.73,0.84,0.94,1.41,3.08,3.12,3.21,4.05,4.06,4.17,6.59,6.63},
  3/{0.89,0.93,1.06,3.1,3.19,3.28,3.89,3.94,4.13,4.2,4.31,6.63,6.65,6.78},
  4/{0.87,0.96,3.09,3.15,3.23,3.73,3.96,4.06,4.19,5.98,6.66,6.68,6.75},
  5/{0.44,0.92,0.97,3.05,3.22,3.27,4.01,4.07,4.17,5.76,6.57,6.73,6.75},
  6/{0.92,1.0,1.73,2.85,3.18,3.19,4.0,4.07,6.53,6.66,6.73},
  7/{0.87,0.92,3.13,3.13,3.27,3.99,4.06,4.17,6.44,6.61,6.72,7.13}}{
  \foreach \s in \ss {\draw[pcBlue, line width=0.7pt] (\s,0.35*\y) -- (\s,0.35*\y+0.25);}}
\draw[arr] (0,-0.3) -- (7.5,-0.3); \node[lbl, anchor=north] at (3.75,-0.32) {时间};
\node[lbl, anchor=east, align=right] at (-0.05,1.35) {八个\\神经元};
\node[lbl, text=pcAmber!70!black, anchor=south] at (3.55,2.95) {集体爆发};
}{放任不管的芯片培养物。单独的咔嗒很少见，而时不时几乎整个网络会一起爆发——然后安静下来，直到连接补足储备。}

\section*{没有多巴胺的老师}

没有多巴胺，怎样教这样的网络？最早的方法是电学的。例如：刺激网络，直到它给出需要的反应——然后立即停止。这个反应就被固定下来了。在另一个实验中，培养物控制一个简单的虚拟“身体”，并学会把它引向目标。在所有这些实验里，老师都是由工程师选定的电流模式。

\section*{靠闪光释放的多巴胺}

有人尝试过把老师做成化学的。在一个培养活神经元小团块的平台上，研究者往营养液里加入了装在光敏“笼子”里的多巴胺，笼子让它无法起作用。每当网络对刺激的反应比以前更强，就“奖励”它：一道紫外线闪光破坏笼子，多巴胺被释放出来。经过二百四十次重复，反应总体上增强了。但作者们坦白地写道，仅凭这样一个实验，无法断定原因就是多巴胺。

\section*{小车上的杆}

在最近的一个实验中，神经元小团块被训练去平衡小车上的一根杆——这是学习程序的一个经典问题。给出什么样的训练脉冲，由一个程序来选择。程序选出的信号比随机信号效果更好，但休息四十五分钟后，学到的东西就消失了。而如果阻断\nt{GLUT}连接，学习就完全无法进行：学到的东西正是存在那里。老师依然在外面——只不过现在不是工程师，而是程序。

\section*{活的储备池}

还有另一种思路：根本不训练活的网络。而是把它当作一个“储备池”——一个丰富而纠缠的系统，由一个简单的电子计数器去学着解读它的反应。一个活的神经元团块就这样帮助分辨过声音。不过，据作者们说，团块内部的连接在工作期间也发生了变化。

\section*{代价与问题}

芯片上的一百万个神经元大约消耗万分之一瓦——远少于它们周围的电子设备。而且这个试验台自己无法决定什么算成功：这个角色总是由外面的某个人来扮演。随着这样的系统越长越大，还冒出了一个至今无人解决的伦理问题。

\fullversion{21}
